\section{Pérdida de la Fuente Radiactiva}

\subsection*{Clasificación: II-B}

\subsubsection{Si la fuente ha sido robada}
Notifique inmediatamente a las oficinas generales y a la C.N.S.N.S.

\subsubsection{Forme brigadas de búsqueda}
Establezca un centro de control de emergencias (CCE) desde donde se puedan coordinar las acciones y centralizar la información.

\subsubsection{Recabe la información necesaria}
Para planificar las labores de búsqueda (planos de la ciudad, personal disponible, vehículos, etc.) de manera sistemática y de tal forma que se opere las 24 horas.

\subsubsection{Coloque avisos para el público}
En los lugares aledaños donde se haya extraviado la fuente radiactiva, previniendo del peligro que representa. Incluya un dibujo o fotografía y advierta de no tocarla, lavarse las manos y avisar de inmediato en caso de no encontrarla.

\subsubsection{Coordine sus acciones con las autoridades civiles}
Y de común acuerdo con la C.N.S.N.S.

\subsubsection{Si la fuente se perdió durante las maniobras de trabajo normal}
Forme brigadas de búsqueda y rastree sistemáticamente el área de trabajo.